\chapter{\ifenglish Introduction\else บทนำ\fi}

\section{\ifenglish Project rationale\else ที่มาและความสำคัญของโครงงาน\fi}
    
    \hspace{1.5em}อาหารพื้นเมืองภาคเหนือของประเทศไทยเป็นส่วนหนึ่งของมรดกทางวัฒนธรรมที่มีเอกลักษณ์ \\ทั้งในด้านรสชาติ วัตถุดิบ และวิธีการปรุงอาหาร 
    อย่างไรก็ตาม ในปัจจุบันพบว่าความรู้เกี่ยวกับอาหารเหล่านี้ในกลุ่มคนรุ่นใหม่ยังมีข้อจำกัด โดยเฉพาะในเชิงลึก เช่น ส่วนผสมและกระบวนการทำอาหาร

    ปัญหาดังกล่าวสะท้อนให้เห็นถึงช่องว่างในการถ่ายทอดองค์ความรู้ด้านวัฒนธรรมอาหารสู่คนรุ่นใหม่ \\โครงงานนี้จึงนำเสนอแนวทางในการใช้ “เกม” เป็นสื่อกลางในการเรียนรู้ โดยผสมผสานความบันเทิงเข้ากับเนื้อหาทางวัฒนธรรม เพื่อให้ผู้เล่นสามารถเรียนรู้ผ่านประสบการณ์การเล่นโดยไม่รู้สึกว่ากำลังศึกษาเนื้อหาเชิงวิชาการ

    เกมที่พัฒนาขึ้นเป็นเกมแนว 2D Roguelike ที่ผสานระบบการต่อสู้ การสำรวจ และการทำอาหารเข้าด้วยกัน โดยมีการออกแบบระบบให้ผู้เล่นต้องเก็บวัตถุดิบจากมอนสเตอร์ นำมาปรุงอาหาร และใช้ผลลัพธ์จากอาหารในการพัฒนาความสามารถของตัวละคร แนวทางนี้ช่วยให้เนื้อหาด้านอาหารถูกเชื่อมโยงเข้ากับกลไกหลักของเกมอย่างเป็นธรรมชาติ

\section{\ifenglish Objectives\else วัตถุประสงค์ของโครงงาน\fi}

    \hspace{1.5em}โครงงานนี้มีวัตถุประสงค์ทั้งในด้านการออกแบบเกม การถ่ายทอดองค์ความรู้ และการประยุกต์ใช้หลักการทางวิศวกรรมคอมพิวเตอร์ โดยสามารถแจกแจงได้ดังนี้

    \begin{enumerate}
    \item เพื่อออกแบบและพัฒนาเกมแนว 2D Roguelike ที่มีระบบการทำอาหารเป็นกลไกหลัก โดยเชื่อมโยงระบบการต่อสู้ การเก็บวัตถุดิบ และการปรุงอาหารเข้าด้วยกันในรูปแบบ gameplay loop ที่ชัดเจน
    \item เพื่อส่งเสริมและถ่ายทอดความรู้เกี่ยวกับอาหารพื้นเมืองภาคเหนือของประเทศไทย โดยนำเสนอข้อมูลเกี่ยวกับวัตถุดิบ รสชาติ และกระบวนการปรุงอาหารในรูปแบบที่เข้าใจง่าย และสอดแทรกอยู่ใน\\กระบวนการเล่นเกม
    \item เพื่อพัฒนาระบบเกมโดยอาศัยหลักการทางวิศวกรรมคอมพิวเตอร์ เช่น การออกแบบสถาปัตยกรรม\\ของระบบ การพัฒนาพฤติกรรมสำหรับตัวละคร และการออกแบบอัลกอริทึมสำหรับระบบให้คะแนน
    \item เพื่อสร้างประสบการณ์ผู้เล่นที่เข้าถึงได้ง่ายและส่งเสริมการเรียนรู้ผ่านการเล่น โดยออกแบบอินเทอร์เฟซและระบบให้เข้าใจง่าย รองรับผู้เล่นที่ไม่มีพื้นฐานทั้งด้านเกมและการทำอาหาร
\end{enumerate}
\newpage
\section{\ifenglish Project scope\else ขอบเขตของโครงงาน\fi}

    \hspace{1.5em}โครงงานนี้มุ่งเน้นการพัฒนาเกมต้นแบบ โดยกำหนดขอบเขตของระบบให้สอดคล้องกับระยะเวลาและ\\ทรัพยากรที่มีอยู่ ทั้งในด้านฮาร์ดแวร์และซอฟต์แวร์ ดังนี้

\subsection{\ifenglish Hardware scope\else ขอบเขตด้านฮาร์ดแวร์\fi}

    \hspace{1.5em}การพัฒนาและทดสอบระบบจะดำเนินการบนอุปกรณ์คอมพิวเตอร์ส่วนบุคคล โดยมีขอบเขตดังนี้

    \begin{itemize}
    \item ใช้คอมพิวเตอร์สำหรับการพัฒนาเกม เช่น การเขียนโปรแกรม การออกแบบระบบ และการทดสอบ
    \item รองรับการรันเกมบนแพลตฟอร์ม Desktop (Windows) เป็นหลัก
    \item ไม่ครอบคลุมการพัฒนาเพื่อแพลตฟอร์มอื่น เช่น Mobile หรือ Console
    \item ไม่รวมถึงการใช้อุปกรณ์เฉพาะทาง เช่น VR/AR หรืออุปกรณ์เสริมพิเศษ
    \end{itemize}

    ขอบเขตดังกล่าวช่วยลดความซับซ้อนของระบบ และทำให้สามารถควบคุมคุณภาพของการพัฒนาได้\\ภายในข้อจำกัดของโครงงาน

\subsection{\ifenglish Software scope\else ขอบเขตด้านซอฟต์แวร์\fi}

    \hspace{1.5em}โครงงานนี้ครอบคลุมการออกแบบและพัฒนาระบบซอฟต์แวร์ของเกมในระดับต้นแบบ โดยมีขอบเขต\\ดังนี้

    \begin{itemize}
        \item พัฒนาเกมในรูปแบบ 2D Roguelike โดยใช้ Game Engine เช่น Unity
        \item ออกแบบและพัฒนาระบบหลักของเกม ได้แก่
            \begin{itemize}
                \item ระบบต่อสู้ (Combat System) 
                \item ระบบวัตถุดิบ (Ingredient System)
                \item ระบบการทำอาหาร (Cooking System)
                \item ระบบแข่งขันทำอาหาร (Competition System)
            \end{itemize}
        \item พัฒนาระบบพฤติกรรมตัวละคร สำหรับ
            \begin{itemize}
                \item มอนสเตอร์ (Monster)
                \item ตัวละครคู่แข่ง (Competitor)
                \item ตัวละครผู้ช่วย (Companion)
            \end{itemize}
        \item ออกแบบระบบให้คะแนนอาหารโดยใช้อัลกอริทึมเชิงคณิตศาสตร์
        \item พัฒนาระบบอินเทอร์เฟซผู้ใช้ (User Interface) ที่เข้าใจง่ายและเหมาะสมกับผู้เล่นทั่วไป
    \end{itemize}

    ทั้งนี้ โครงงานจะเน้นการพัฒนาในระดับ functional prototype โดยอาจยังไม่ครอบคลุมการปรับปรุงในระดับ production เช่น optimization ขั้นสูง หรือเนื้อหาที่สมบูรณ์ทั้งหมด

\section{\ifenglish Expected outcomes\else ประโยชน์ที่ได้รับ\fi}

    \hspace{1.5em}โครงงานนี้ก่อให้เกิดประโยชน์ในหลายด้าน ทั้งในเชิงวิชาการ เทคโนโลยี และสังคม ดังนี้

    \begin{enumerate}
        \item \textbf{ด้านการเรียนรู้และวัฒนธรรม}\newline
        \hspace{1.5em}ผู้เล่นสามารถเรียนรู้เกี่ยวกับอาหารพื้นเมืองภาคเหนือ ทั้งในด้านวัตถุดิบ รสชาติ และกระบวนการปรุงอาหาร ผ่านกลไกของเกม ซึ่งช่วยส่งเสริมการอนุรักษ์และเผยแพร่วัฒนธรรมในรูปแบบที่เข้าถึงคนรุ่นใหม่ได้ง่าย
        \item \textbf{ด้านประสบการณ์ผู้เล่น}\newline
        \hspace{1.5em}เกมที่พัฒนาขึ้นช่วยสร้างประสบการณ์การเรียนรู้แบบมีส่วนร่วม ทำให้ผู้เล่นสามารถเรียนรู้ผ่านการ\\ลงมือทำ ส่งผลให้เกิดความเข้าใจและการจดจำที่ดียิ่งขึ้น
        \item \textbf{ด้านวิศวกรรมคอมพิวเตอร์}\newline
        \hspace{1.5em}ผู้พัฒนาได้ประยุกต์ใช้ความรู้ด้านวิศวกรรมคอมพิวเตอร์ เช่น การออกแบบสถาปัตยกรรมระบบ \\การพัฒนาพฤติกรรมหลายรูปแบบ และการออกแบบอัลกอริทึมสำหรับระบบเกม ซึ่งสามารถนำไป\\ต่อยอดในงานด้านการพัฒนาเกมหรือซอฟต์แวร์อื่น ๆ ได้
        \item \textbf{ด้านการออกแบบระบบเกม}\newline
        \hspace{1.5em}โครงงานนี้แสดงให้เห็นถึงการออกแบบระบบเกมที่เชื่อมโยงหลายองค์ประกอบเข้าด้วยกัน เช่น ระบบการต่อสู้ ระบบการทำอาหาร และระบบความก้าวหน้าของตัวละคร ซึ่งเป็นแนวทางที่สามารถนำไปพัฒนาเกมในรูปแบบอื่นได้
        \item \textbf{ด้านการพัฒนาทักษะของผู้พัฒนา}\newline
        \hspace{1.5em}ผู้พัฒนาได้ฝึกทักษะในหลายด้าน เช่น การเขียนโปรแกรม การออกแบบระบบ การทำงานเป็นทีม \\และการแก้ปัญหาเชิงวิศวกรรม ซึ่งเป็นทักษะสำคัญในการทำงานจริงในสายอุตสาหกรรม
    \end{enumerate}

\section{\ifenglish Technology and tools\else เทคโนโลยีและเครื่องมือที่ใช้\fi}

    \hspace{1.5em}โครงงานนี้ใช้เทคโนโลยีและเครื่องมือทั้งด้านฮาร์ดแวร์และซอฟต์แวร์ เพื่อสนับสนุนการออกแบบ พัฒนา และทดสอบระบบเกมให้มีประสิทธิภาพ โดยสามารถแบ่งออกได้ดังนี้

\subsection{\ifenglish Hardware technology\else เทคโนโลยีด้านฮาร์ดแวร์\fi}

    \hspace{1.5em}การพัฒนาเกมในโครงงานนี้ดำเนินการบนคอมพิวเตอร์ส่วนบุคคล ซึ่งใช้สำหรับการเขียนโปรแกรม \\ออกแบบระบบ และทดสอบการทำงานของเกม โดยมีองค์ประกอบหลักดังนี้
    
    \begin{itemize}
        \item หน่วยประมวลผลกลาง (CPU) สำหรับการประมวลผลคำสั่งและการทำงานของระบบ
        \item หน่วยความจำหลัก (RAM) สำหรับรองรับการทำงานของโปรแกรมและ Game Engine
        \item หน่วยประมวลผลกราฟิก (GPU) สำหรับแสดงผลกราฟิกของเกม
        \item อุปกรณ์จัดเก็บข้อมูล (Storage) สำหรับเก็บไฟล์โปรเจกต์และทรัพยากรของเกม
    \end{itemize}

    ฮาร์ดแวร์ที่ใช้เป็นอุปกรณ์ทั่วไปที่สามารถรองรับการพัฒนาเกม 2D ได้ โดยไม่จำเป็นต้องใช้อุปกรณ์\\เฉพาะทาง

\subsection{\ifenglish Software technology\else เทคโนโลยีด้านซอฟต์แวร์\fi}

    \hspace{1.5em}โครงงานนี้ใช้เครื่องมือและเทคโนโลยีด้านซอฟต์แวร์ในการพัฒนาเกม ดังนี้

    \begin{itemize}
        \item Game Engine: Unity\newline
        \hspace{1.5em}ใช้สำหรับพัฒนาเกม 2D รวมถึงการจัดการระบบต่าง ๆ เช่น physics, animation และ \\scene management
        \item Programming Language: C\#\newline
        \hspace{1.5em}ใช้ในการพัฒนา logic ของเกม เช่น ระบบต่อสู้ ระบบ AI และระบบการทำอาหาร
        \item Graphic Design Tools
            \begin{itemize}
                \item \textbf{Blender} ใช้สำหรับการสร้างและปรับแต่งโมเดลหรือองค์ประกอบกราฟิกที่เกี่ยวข้อง
                \item \textbf{Clip Studio Paint} ใช้สำหรับออกแบบภาพ 2D เช่น ตัวละคร ไอคอน และองค์ประกอบของอินเทอร์เฟซ
            \end{itemize}
        \item Graphic and Audio Assets\newline
        \hspace{1.5em}ใช้ทรัพยากรภาพและเสียงจากแหล่งฟรี และบางส่วนสร้างขึ้นเองเพื่อให้เหมาะสมกับธีมของเกม
        \item Version Control\newline
        \hspace{1.5em}เช่น Git สำหรับจัดการเวอร์ชันของโค้ดและการทำงานร่วมกันภายในทีม
    \end{itemize}

    เทคโนโลยีที่เลือกใช้ทั้งหมดมุ่งเน้นความเหมาะสมกับขอบเขตของโครงงาน และสามารถรองรับการ\\พัฒนาเกมต้นแบบได้อย่างมีประสิทธิภาพ

\section{\ifenglish Project plan\else แผนการดำเนินงาน\fi}

\begin{plan}{11}{2025}{3}{2026}
    \planitem{11}{2025}{11}{2025}{กำหนดหัวข้อและวัตถุประสงค์ของโครงงาน}
    \planitem{12}{2025}{12}{2025}{ศึกษาข้อมูลด้านอาหารและวัฒนธรรมภาคเหนือ}
    \planitem{12}{2025}{1}{2026}{ออกแบบระบบของเกมในภาพรวม}
    \planitem{1}{2026}{3}{2026}{ออกแบบเชิงลึกของระบบ}
    \planitem{1}{2026}{3}{2026}{ออกแบบ Gameplay และประสบการณ์ผู้เล่น}
    \planitem{3}{2026}{3}{2026}{การเก็บข้อมูลจากแบบสอบถาม}
\end{plan}
\newpage
\section{\ifenglish Roles and responsibilities\else บทบาทและความรับผิดชอบ\fi}

    \hspace{1.5em}โครงงานนี้ดำเนินงานโดยทีมพัฒนา ซึ่งมีการแบ่งบทบาทและความรับผิดชอบอย่างชัดเจน โดยแต่ละสมาชิกมีหน้าที่เฉพาะด้าน พร้อมทั้งมีส่วนร่วมในการพัฒนาระบบโดยรวม ดังนี้

    \begin{itemize}
        \item \textbf{ชลกร สุทธเวช (660610746)}\newline
        \hspace{1.5em}รับผิดชอบด้านการพัฒนาระบบการทำอาหาร ในเชิงกลไกของเกม เช่น การจัดการวัตถุดิบ \\กระบวนการปรุงอาหาร และการเชื่อมโยงระบบการทำอาหารเข้ากับระบบสกิลและการต่อสู้ \\เพื่อให้เกิดความสัมพันธ์ระหว่างระบบต่าง ๆ ภายในเกม
        \item \textbf{ณัชชา คำปวง (660610751)}\newline
        \hspace{1.5em}รับผิดชอบด้านการออกแบบและพัฒนา Game Assets โดยใช้เครื่องมือ เช่น Blender ในการสร้าง\\ตัวละคร ไอคอน และองค์ประกอบของอินเทอร์เฟซ รวมถึงมีส่วนในการพัฒนาระบบการทำอาหาร \\โดยเน้นด้านการออกแบบประสบการณ์ผู้ใช้ (User Interface / User Experience) ให้เข้าใจง่ายและใช้งานสะดวก
        \item \textbf{บุรัสกร ศรีขาว (660612150)}\newline
        \hspace{1.5em}รับผิดชอบด้านการออกแบบสถาปัตยกรรมของระบบเกม และพัฒนาในส่วนของระบบการต่อสู้ \\โดยครอบคลุมการออกแบบพฤติกรรมของมอนสเตอร์ การจัดการรูปแบบการโจมตี 
        และการตอบสนองของผู้เล่น รวมถึงมีส่วนร่วมในการพัฒนาองค์ประกอบด้านกราฟิกภายในเกม นอกจากนี้ยังรับผิดชอบในการปรับสมดุลของระบบการต่อสู้ เพื่อให้เกิดประสบการณ์การเล่นที่เหมาะสม
    \end{itemize}

    นอกจากนี้ สมาชิกทุกคนในทีมมีส่วนร่วมในการทดสอบระบบ (Testing) การแก้ไขข้อผิดพลาด \\(Debugging) และการปรับปรุงประสิทธิภาพของเกมอย่างต่อเนื่อง รวมถึงมีการทำงานร่วมกันในการ\\ออกแบบและปรับปรุงระบบเพื่อให้ทุกองค์ประกอบสามารถทำงานสอดคล้องกันได้อย่างมีประสิทธิภาพ

\section{\ifenglish%
Impacts of this project on society, health, safety, legal, and cultural issues
\else%
ผลกระทบด้านสังคม สุขภาพ ความปลอดภัย กฎหมาย และวัฒนธรรม
\fi}

    \hspace{1.5em}การพัฒนาเกมในโครงงานนี้อาจก่อให้เกิดผลกระทบในหลายด้าน ทั้งเชิงบวกและข้อควรระวัง \\โดยสามารถพิจารณาได้ดังนี้

\subsection{\ifenglish Social impact\else ด้านสังคม\fi}
    \hspace{1.5em}โครงงานนี้มีส่วนช่วยส่งเสริมการเรียนรู้และการเข้าถึงองค์ความรู้ด้านอาหารพื้นเมืองภาคเหนือของ\\ประเทศไทยผ่านสื่อเกม ซึ่งเป็นรูปแบบที่เข้าถึงคนรุ่นใหม่ได้ง่าย ส่งผลให้เกิดการเผยแพร่วัฒนธรรมในวงกว้าง 
    \\อย่างไรก็ตาม การเล่นเกมเป็นเวลานานอาจส่งผลต่อพฤติกรรมการใช้เวลา ผู้พัฒนาจึงควรออกแบบเกมให้มีความเหมาะสมและไม่ส่งเสริมการเล่นอย่างต่อเนื่องเกินไป

\subsection{\ifenglish Health impact\else ด้านสุขภาพ\fi}
    \hspace{1.5em}เกมที่พัฒนาขึ้นสามารถสร้างความเพลิดเพลินและช่วยผ่อนคลายความเครียดให้กับผู้เล่นได้
    \\ในทางกลับกัน หากผู้เล่นใช้เวลาเล่นเกมมากเกินไป อาจส่งผลกระทบต่อสุขภาพ เช่น อาการเมื่อยล้าทางสายตา หรือการขาดการเคลื่อนไหว ดังนั้นควรมีการออกแบบให้เหมาะสม เช่น การแบ่งช่วงการเล่นเป็นรอบ

\subsection{\ifenglish Safety impact\else ด้านความปลอดภัย\fi}
    \hspace{1.5em}เนื้อหาของเกมถูกออกแบบให้อยู่ในระดับที่เหมาะสมกับผู้เล่นทั่วไป โดยไม่มีเนื้อหาที่รุนแรงหรือ\\ไม่เหมาะสม นอกจากนี้ เกมไม่ได้มีการเก็บข้อมูลส่วนบุคคลของผู้เล่นในระดับที่ละเอียด จึงช่วยลดความเสี่ยงด้านความปลอดภัยของข้อมูล

\subsection{\ifenglish Legal impact\else ด้านกฎหมาย\fi}
    \hspace{1.5em}การพัฒนาเกมในโครงงานนี้คำนึงถึงกฎหมายที่เกี่ยวข้อง เช่น
    \begin{itemize}
        \item การใช้ทรัพยากร (assets) ที่ไม่มีการละเมิดลิขสิทธิ์
        \item การใช้ซอฟต์แวร์และเครื่องมือที่มีสิทธิ์ใช้งานถูกต้อง
    \end{itemize}
    \hspace{1.5em}นอกจากนี้ เนื้อหาของเกมไม่มีการละเมิดหรือบิดเบือนข้อมูลทางวัฒนธรรมในลักษณะที่ก่อให้เกิดความเข้าใจผิด

\subsection{\ifenglish Cultural impact\else ด้านวัฒนธรรม\fi}
    \hspace{1.5em}โครงงานนี้มีบทบาทในการส่งเสริมและเผยแพร่วัฒนธรรมอาหารภาคเหนือของประเทศไทย โดยนำเสนอในรูปแบบที่เข้าใจง่ายและน่าสนใจ\par
    อย่างไรก็ตาม ผู้พัฒนาตระหนักถึงความสำคัญของการนำเสนอข้อมูลทางวัฒนธรรมอย่างถูกต้องและ\\เหมาะสม เพื่อหลีกเลี่ยงการบิดเบือนหรือทำให้เกิดความเข้าใจคลาดเคลื่อน\par
    โดยสรุป โครงงานนี้มีผลกระทบเชิงบวกในด้านการเรียนรู้และการอนุรักษ์วัฒนธรรมเป็นหลัก \\พร้อมทั้งมีการคำนึงถึงข้อจำกัดและความเสี่ยงในด้านต่าง ๆ เพื่อให้การพัฒนาเป็นไปอย่างมีความรับผิดชอบต่อสังคม

\section{\ifenglish Budget\else งบประมาณ\fi}

    \hspace{1.5em}โครงงานนี้เป็นการพัฒนาซอฟต์แวร์เป็นหลัก จึงมีค่าใช้จ่ายโดยรวมในระดับต่ำ โดยเน้นการใช้เครื่องมือฟรีและทรัพยากรที่มีอยู่แล้ว อย่างไรก็ตาม งบประมาณของโครงงานจะถูกใช้ในส่วนของทรัพยากรดิจิทัลและเครื่องมือสนับสนุนการพัฒนาเป็นหลัก โดยสามารถสรุปได้ดังนี้

\subsection{\ifenglish Digital Assets\else ทรัพยากรดิจิทัล (Assets)\fi}
    \hspace{1.5em}งบประมาณส่วนใหญ่ของโครงงานจะถูกใช้กับการจัดหา Asset ที่จำเป็นสำหรับการพัฒนาเกม เช่น
    \begin{itemize}
        \item กราฟิก (Sprites, UI, Effects)
        \item โมเดลหรือองค์ประกอบภาพเพิ่มเติม
        \item เสียงและดนตรีประกอบ (Sound Effects, Background Music)
    \end{itemize}
    \hspace{1.5em}การเลือกใช้ Asset จะพิจารณาจากความเหมาะสมกับธีมของเกม และความคุ้มค่า โดยอาจใช้ทั้ง\\ทรัพยากรฟรีและทรัพยากรที่มีค่าใช้จ่าย

\subsection{\ifenglish Development Tools\else เครื่องมือและตัวช่วยในการพัฒนา\fi}
    \hspace{1.5em}อาจมีค่าใช้จ่ายในส่วนของเครื่องมือหรือบริการที่ช่วยเพิ่มประสิทธิภาพในการพัฒนา เช่น
    \begin{itemize}
        \item ปลั๊กอินเสริมสำหรับ Game Engine
        \item เครื่องมือช่วยออกแบบหรือจัดการ Asset
        \item บริการออนไลน์ที่สนับสนุนการพัฒนา
    \end{itemize}
    \hspace{1.5em}อย่างไรก็ตาม โครงงานนี้พยายามลดค่าใช้จ่ายในส่วนนี้โดยเลือกใช้เครื่องมือที่ไม่มีค่าใช้จ่ายเป็นหลัก

\subsection{\ifenglish Software and Hardware Costs\else ค่าใช้จ่ายด้านซอฟต์แวร์และฮาร์ดแวร์\fi}
    \begin{itemize}
        \item ใช้ Game Engine (Unity) และเครื่องมือหลักอื่น ๆ ในเวอร์ชันฟรี
        \item ใช้คอมพิวเตอร์ส่วนบุคคลของผู้พัฒนา จึงไม่มีค่าใช้จ่ายเพิ่มเติมด้านฮาร์ดแวร์
    \end{itemize}

\subsection{\ifenglish Summary of Budget\else สรุปงบประมาณ\fi}
    \hspace{1.5em}โดยรวมแล้ว งบประมาณของโครงงานจะถูกใช้ในส่วนของ Asset และเครื่องมือสนับสนุนเป็นหลัก \\ในขณะที่ค่าใช้จ่ายด้านซอฟต์แวร์และฮาร์ดแวร์อยู่ในระดับต่ำ ทำให้สามารถควบคุมงบประมาณของโครงงานได้อย่างมีประสิทธิภาพ